\documentclass{article}

% ------------------------------------------------------------
% arxiv.sty is in this directory (downloaded from kourgeorge/arxiv-style)
% Build:  pdflatex main.tex && bibtex main && pdflatex main.tex && pdflatex main.tex
% ------------------------------------------------------------
\usepackage{arxiv}

\usepackage[utf8]{inputenc}
\usepackage[T2A,T1]{fontenc}    % T2A enables Cyrillic glyphs
\usepackage[russian,english]{babel}
\usepackage{hyperref}
\usepackage{url}
\usepackage{booktabs}
\usepackage{amsfonts}
\usepackage{amsmath}
\usepackage{nicefrac}
\usepackage{microtype}
\usepackage{graphicx}
\usepackage{subcaption}
\usepackage{xcolor}

\title{
  Bөlinbeitin\footnote{Kazakh \emph{бөлінбейтін} — ``indivisible''; refers to the morpheme,
  the smallest meaning-bearing unit of grammar.}: \\
  A Morphological Alignment Audit of Subword Tokenization \\
  for Agglutinative Kazakh and Fusional Russian
}

\author{
  Altair Zhambyl \\
  Independent Researcher \\
  Almaty, Kazakhstan \\
  \texttt{hey@evgenyastapov.com}
}

% ------------------------------------------------------------
\begin{document}
\maketitle

\begin{abstract}
% ABSTRACT WRITING NOTES (delete before submission):
% 150-250 words. Write LAST, after Results.
% Structure: (1) one-sentence context (2) gap (3) what you did (4) main finding (5) implication.
% Numbers go in the abstract when they are surprising.
Subword tokenization with byte-pair encoding (BPE) has become the de-facto preprocessing
step for large language models, yet its behaviour on morphologically rich languages
remains under-studied. We present a controlled morphological alignment audit of six
production tokenizers (GPT-4 \texttt{cl100k\_base}, LLaMA-3, mT5, BLOOM, Cohere Aya,
and Qwen) on a native-annotated test set of Kazakh, an agglutinative Turkic language;
Russian, a fusional Indo-European language; and English, an analytic baseline.
We introduce two new metrics --- \emph{Morphological Boundary Alignment} (MBA) and
\emph{Suffix Integrity Score} (SIS) --- that quantify how often token boundaries
respect morpheme boundaries and how often inflectional suffixes survive as single
tokens. \textsc{TODO Day 5: main quantitative finding goes here, e.g.\
``We find that Kazakh requires \(X\times\) more tokens than English to encode equivalent
content, and that fewer than \(Y\%\) of token boundaries align with morpheme boundaries
even for the multilingual mT5 tokenizer.''} We release the annotated test sets under
CC-BY-4.0 to support future work on morphology-aware tokenization for Turkic and
other agglutinative language families.
\end{abstract}

% ============================================================
\section{Introduction}
\label{sec:intro}

% WRITING NOTES (Day 6):
% 1-1.5 pages. Funnel structure:
% (a) Hook: general problem
% (b) Specific gap in literature
% (c) Contributions (bullet list)
% (d) Roadmap

% Hook
\textsc{TODO Day 6: opening hook.}
Large language models depend on subword tokenizers to bridge raw text and the
fixed-vocabulary input layer of the network~\citep{sennrich2016bpe}.
Tokenization choices shape both the cost of training and the inductive biases
the model brings to morphology, syntax, and cross-lingual generalization.
% Gap
Recent work has documented dramatic per-token cost disparities across
languages~\citep{petrov2023tokenizer, ahia2023languages},
but the \emph{linguistic structure} underlying those disparities --- particularly
the relationship between token boundaries and morpheme boundaries in
morphologically rich, low-resource languages such as Kazakh --- remains largely
unmeasured.

% Contributions
We make the following contributions:
\begin{itemize}
  \item A native-annotated morphological test set of \(500\) Kazakh words,
        \(500\) Russian words, and \(200\) English words, released under
        CC-BY-4.0.\footnote{\url{TODO github url}}
  \item Two novel metrics for measuring token--morpheme structural alignment:
        \emph{Morphological Boundary Alignment} (MBA, \S\ref{sec:metrics}) and
        \emph{Suffix Integrity Score} (SIS, \S\ref{sec:metrics}).
  \item A side-by-side audit of six production tokenizers across the three
        languages, with all per-token outputs released for reproducibility.
  \item A list of minimum design requirements for a morphology-aware tokenizer
        targeting agglutinative Turkic languages (\S\ref{sec:discussion}).
\end{itemize}

The remainder of the paper is organized as follows.
Section~\ref{sec:related} reviews prior work.
Section~\ref{sec:method} describes our test sets, the tokenizers under audit, and our metrics.
Section~\ref{sec:results} presents the audit results.
Section~\ref{sec:discussion} discusses design implications.
Section~\ref{sec:limitations} acknowledges limitations.
Section~\ref{sec:conclusion} concludes.

% ============================================================
\section{Related Work}
\label{sec:related}

% WRITING NOTES (Day 4):
% Three short paragraphs. ~0.7 pages total.

\paragraph{Subword tokenization.}
\textsc{TODO Day 4.} BPE~\citep{sennrich2016bpe}, WordPiece, Unigram LM.
End with a sentence connecting to why this matters cross-lingually.

\paragraph{Cross-lingual tokenization fairness.}
\textsc{TODO Day 4.} Petrov et al.~\citep{petrov2023tokenizer} showed that token
counts vary by a factor of up to 15 across the 22 languages they audited,
penalizing speakers of low-resource languages through both higher inference cost
and reduced effective context window. Ahia et al.~\citep{ahia2023languages}
extended this analysis to API pricing. Our work goes one level deeper: instead of
asking \emph{how many} tokens a language consumes, we ask \emph{whether those tokens
correspond to linguistically meaningful units}.

\paragraph{Morphological tokenization for agglutinative languages.}
\textsc{TODO Day 4.} Toraman et al.~\citep{toraman2023turkish} audited tokenizers
on Turkish, an agglutinative language closely related to Kazakh, and found
substantial morphology--tokenization mismatch. Bostrom and
Durrett~\citep{bostrom2020bpe} argued BPE is suboptimal for language model
pretraining. We contribute the first such audit specifically for Kazakh, the
fourth-most-spoken Turkic language with $\sim$13M speakers but limited NLP coverage.

% ============================================================
\section{Method}
\label{sec:method}

% WRITING NOTES (Day 2-3):
% Be precise enough that a stranger can reproduce your numbers.

\subsection{Test Sets}
\label{sec:testsets}
We construct three parallel test sets of single tokens (words) with morphological
annotation:
\begin{itemize}
  \item \textbf{Kazakh} (500 words): collected and annotated by the author,
        a native speaker. Each entry records the orthographic word, a hyphen-separated
        morpheme analysis, a Leipzig-style gloss~\citep{lehmann2004leipzig}, part-of-speech,
        and a test-case category.
  \item \textbf{Russian} (500 words): annotated similarly; cross-checked against the
        \texttt{pymorphy3} morphological analyzer.
  \item \textbf{English} (200 words): a smaller baseline reflecting English's
        comparatively poor morphological inventory.
\end{itemize}
Test cases are distributed across categories including bare nouns, case-marked nouns,
possessive constructions, plural marking, verbal inflection, negation, and derivational
morphology. The Kazakh set additionally targets the seven case markers (\emph{septik})
and the six possessive paradigms (\emph{tauelldik}). \textsc{TODO Day 3: insert
distribution table.}

\subsection{Tokenizers Under Audit}
\label{sec:tokenizers}
We audit six production tokenizers, chosen for their use in widely deployed LLMs and
their declared training-data multilinguality. \textsc{TODO Day 2: insert table from
TOPIC.md with vocab sizes, model families, and language coverage notes.}

\subsection{Metrics}
\label{sec:metrics}

\paragraph{Tokens per word (TPW) and per character (TPC).}
Standard efficiency measures: $\textrm{TPW}(t, w) = |t(w)|$ where $t(\cdot)$ is the
tokenizer applied to word $w$, and $\textrm{TPC} = \textrm{TPW}/|w|$ averaged over
the test set.

\paragraph{Morphological Boundary Alignment (MBA).}
Given a word $w$ with morpheme boundaries $B_m = \{b_1, \dots, b_k\}$ (character offsets
between morphemes) and tokenizer output with token boundaries $B_t$, we define
\[
  \textrm{MBA}(w, t) \;=\; \frac{|B_t \cap B_m|}{|B_t|}.
\]
MBA $= 1$ means every token boundary the tokenizer introduced respects a real morpheme
boundary; MBA $= 0$ means every boundary cuts a morpheme in half.

\paragraph{Suffix Integrity Score (SIS).}
For each inflectional suffix $s$ in the annotated test set, $\textrm{SIS}(s, t) = 1$
if $s$ appears as a single contiguous substring of one token in $t(w)$, else $0$.
We report mean SIS per suffix category (case markers, possessive markers, person markers).

\subsection{Reproducibility Protocol}
All tokenizer outputs, raw metric tables, and the annotated test sets are released
under CC-BY-4.0. A single command (\texttt{python run\_audit.py}) regenerates every
table and figure in the paper.

% ============================================================
\section{Results}
\label{sec:results}

% WRITING NOTES (Day 4-5):
% One subsection per RQ. Lead with a figure. Then prose.

\subsection{RQ1: How many tokens per word?}
\label{sec:rq1}
\textsc{TODO Day 3-4: lead with bar chart Figure.}

\begin{figure}[h]
\centering
\includegraphics[width=0.85\linewidth]{figures/tpw_by_lang_tokenizer.pdf}
\caption{Mean tokens per word across three languages and six tokenizers.
The ``Kazakh tax'' is visible as the consistent leftmost-bar elevation across all tokenizers.}
\label{fig:tpw}
\end{figure}

\textsc{TODO Day 5: 2-3 paragraphs interpreting Figure~\ref{fig:tpw}.}

\subsection{RQ2: How well do token boundaries align with morpheme boundaries?}
\label{sec:rq2}
\textsc{TODO Day 4-5.}

\subsection{RQ3: Do inflectional suffixes survive tokenization?}
\label{sec:rq3}
\textsc{TODO Day 4-5.}

% ============================================================
\section{Discussion: Design Implications}
\label{sec:discussion}

\textsc{TODO Day 5.}
Based on our findings, a morphology-aware tokenizer for agglutinative Turkic
languages should:
\begin{enumerate}
  \item \textsc{TODO 4-6 concrete design requirements.}
\end{enumerate}

% ============================================================
\section{Limitations}
\label{sec:limitations}

% Be honest. Reviewers respect this.
\begin{itemize}
  \item Test sets were annotated by a single native speaker; future work should
        include inter-annotator agreement studies.
  \item The Kazakh test set is small (500 words) relative to a real corpus,
        and was not stratified by frequency.
  \item We measure structural alignment, not downstream task impact;
        Paper 2 in this research programme will close that gap.
  \item We cover only the standard literary Kazakh orthography (Cyrillic script,
        2024 Latin reform not included) and do not address dialectal variation.
  \item Six tokenizers is a snapshot; the production-tokenizer landscape moves quickly.
\end{itemize}

% ============================================================
\section{Conclusion and Future Work}
\label{sec:conclusion}

\textsc{TODO Day 6.}
This paper is the first step of a longer research programme on
Kazakh-first language modelling. Immediate next steps include:
(i) training a morphology-informed BPE tokenizer for Kazakh and re-running
this audit on it; (ii) measuring downstream impact on Kazakh NLP tasks; and
(iii) extending the test set to additional Turkic languages
(Kyrgyz, Uzbek, Tatar) for cross-family comparison.

% ============================================================
\section*{Acknowledgements}
The author thanks the open-source community --- particularly the maintainers of
HuggingFace \texttt{tokenizers}, \texttt{tiktoken}, \texttt{sentencepiece},
\texttt{pymorphy3}, and the Apertium project --- whose tools made this audit
possible on a single laptop. \textsc{TODO Day 7: add specific human acknowledgements.}

% ============================================================
\bibliographystyle{plainnat}
\bibliography{references}

\end{document}
